% Source: mathematics/periodic_point_trace_comparison.md
% Source SHA256: 8a7f60d5a7aa55a69094d31df7a8d685d98b23bb5b7790e046ad55d8731d0428
\hypertarget{a-return-trace-comparison-for-general-smooth-phase-space-symbols}{%
\section{A return-trace comparison for general smooth phase-space symbols}\label{a-return-trace-comparison-for-general-smooth-phase-space-symbols}}
\par\smallskip\noindent\emph{Numbering convention: section and theorem references within this chapter use the source module\textquotesingle s local numbering; source Section 1 is subsection \thesection.1. Equation references carry the prefix Ret.}\par\smallskip

\textbf{Status.} Project proof with exact finite algebra checks and separately
labelled floating physical diagnostics. Limited independent internal
analytic checks cover Sections 1--5 and 8--10; they are not external referee
acceptance. The coupling and joint-law extension have been checked by the author.
Priority verification and Lean formalization are outstanding. Base revision:
9ed93cd1bc5e3362c1dbf520f0f391478dfb6969.

This note removes the position-only restriction for a one-period comparison.
Sections 7--9 now identify its Gaussian limit and extend the leading
asymptotic to every replica count at the critical scale. This does not prove
microscopic BGS or random-matrix correlations. The classical kick strength
tends to zero.

\hypertarget{quantization-and-statement}{%
\subsection{Quantization and statement}\label{quantization-and-statement}}

Retain
\[
 \begin{gathered}
 A=\begin{pmatrix}3&4\\2&3\end{pmatrix},\quad
 \lambda=3+2\sqrt2,\quad K=2m,\quad m\ge3\ \hbox{odd},\quad N=b_m,
 \\ b_0=0,\quad b_1=1,\quad b_{j+1}=6b_j-b_{j-1}.
 \end{gathered}
\]
Normalize the metaplectic operator by \(U_0^K=I\).
For integer frequency vectors use the balanced Weyl convention
\[
 \mathcal W_N(v)=(-1)^{v_1v_2}W_N(v),\qquad
 \mathcal W_N(v)\mathcal W_N(w)
       =e^{\pi i\omega(v,w)/N}\mathcal W_N(v+w).             \tag{Ret.P1}
\]
\(W_N\) is the modular-centered convention of the preceding notes.
Unlike \(W_N\), \(\mathcal W_N\) has signs under shifts by \(N\).
In particular \(\mathcal W_N(2Nv)=I\).
Since \(A\equiv I\pmod2\), exact Egorov holds in this convention too.

For a real smooth symbol on \(\mathbb T^2\), write
\[
 f(z)=\sum_{v\in\mathbb Z^2}\widehat f_v e^{2\pi i\langle v,z\rangle},
 \quad F_j=\sum_v\|v\|_1^j|\widehat f_v|,\quad
 H_N=\operatorname{Op}_N(f)=\sum_v\widehat f_v\mathcal W_N(v),
 \quad U_{\epsilon,f}=e^{i\epsilon H_N}U_0 .
\]
In \(F_0\) the weight is one, including at zero frequency.
\(H_N\) is self-adjoint. No parity or mean-zero assumption is needed
for the full trace. If \(f(-z)=f(z)\), the odd block is invariant.

Let
\[
 \mathcal P_K=\{z:(A^{\mathsf T})^Kz=z\pmod{\mathbb Z^2}\},\qquad
 Z_{K,f}(\epsilon)=\frac1{|\mathcal P_K|}
       \sum_{z\in\mathcal P_K}
       \exp\left(i\epsilon\sum_{j=0}^{K-1}f((A^{\mathsf T})^jz)\right).
                                                               \tag{Ret.P2}
\]
All fixed points are counted, including those of smaller primitive period,
with equal weight. This finite group is explicitly
\[
 \mathcal P_K=
 \left\{\left(\frac{u}{8N},\frac{v}{4N}\right):
             0\le u<8N,\ 0\le v<4N\right\},\qquad
                         |\mathcal P_K|=32N^2.             \tag{Ret.P3}
\]

\textbf{Theorem 1 (finite Fourier support).}
Suppose \(\widehat f_v=0\) for \(\|v\|_1>M\), \(M\ge1\).
Put \(\delta=|\epsilon|\), \(u=K\delta F_0\),
\(L=\lceil K/(64M)\rceil\), and
\(E_L(u)=\sum_{\ell\ge L}u^\ell/\ell!\).
For every real \(\epsilon\),
\[
 \left|\frac{\operatorname{Tr}U_{\epsilon,f}^K}{N}
                                  -Z_{K,f}(\epsilon)\right|
 \le \frac{2\pi e^u}{N}
       \left[K\delta F_2+(K\delta F_1)^2+K\delta^2F_1^2\right]
                                      +E_L(u).             \tag{Ret.P4}
\]
For an even symbol the normalized odd trace satisfies the same bound
with the additional term \(e^u/(N-1)\).

\textbf{Theorem 2 (arbitrary fixed smooth symbol).}
For every real \(f\in C^\infty(\mathbb T^2)\), every \(T<\infty\), and \(q>0\),
\[
 \sup_{|t|\le T}
 \left|\frac{\operatorname{Tr}U_{t/\sqrt K,f}^K}{N}
                         -Z_{K,f}(t/\sqrt K)\right|
                        =O_{f,T,q}(K^{-q}).                \tag{Ret.P5}
\]
The same holds for the odd trace when \(f\) is even.
This is superpolynomial accuracy in \(K\asymp\log N\), not in \(N\).
For a fixed Fourier polynomial (Ret.P4) is stronger at this scale:
the error is \(N^{-1+o(1)}\).
It also gives an exponentially small error in \(K\) for a sufficiently
small fixed \(\delta>0\), depending on \(M,F_0\).
These statements concern one arithmetic period \(K\).

\hypertarget{the-first-order-lattice-and-its-32-classes}{%
\subsection{The first-order lattice and its 32 classes}\label{the-first-order-lattice-and-its-32-classes}}

On cyclic vector sequences define
\[
 (\mathcal Dp)_j=p_j-Ap_{j+1},\qquad
 \mathcal C_K=\left\{c\in(\mathbb Z^2)^K:
                         \sum_jA^jc_j=0\pmod N\right\}.    \tag{Ret.P6}
\]
The total-frequency map onto \((\mathbb Z/N\mathbb Z)^2\) is surjective,
so \(\mathcal C_K\) has index \(N^2\).
Also \(\mathcal D(\mathbb Z^{2K})\subset\mathcal C_K\), and
\[
 |\det\mathcal D|=|\det(I-A^K)|=32N^2.
\]
Consequently
\[
 |\mathcal C_K/\mathcal D(\mathbb Z^{2K})|=32,\qquad
       p=\mathcal D^{-1}c\in(1/32)\mathbb Z^{2K}.            \tag{Ret.P7}
\]
The last assertion follows because a finite group of order 32 is
annihilated by 32. Individual classes need not be classified.

There is a uniform cyclic inverse bound
\[
                            \|p\|_1\le2\|\mathcal Dp\|_1.  \tag{Ret.P8}
\]
Let \(P_+,P_-\) be the spectral projectors of \(A\) with eigenvalues
\(\lambda,\lambda^{-1}\). Extend the cyclic forcing periodically.
The unstable inverse is \(-\sum_{h\ge1}\lambda^{-h}P_+c_{j-h}\);
the stable inverse is \(\sum_{h\ge0}\lambda^{-h}P_-c_{j+h}\).
These absolutely convergent sums give the unique cyclic inverse.
Their induced \(\ell^1\) norm is at most
\[
 \frac{\|P_+\|_1}{\lambda-1}
 +\frac{\|P_-\|_1}{1-\lambda^{-1}}
                      =1+\frac1{\sqrt2}<2.
\]
If \(p\) is not integer, none of its vector coordinates is an integer
vector: \(p_j-Ap_{j+1}\in\mathbb Z^2\) and the integral inverse of \(A\)
propagate integrality in both directions.
By (Ret.P7), every coordinate then has norm at least \(1/32\).
Thus every noninteger class satisfies
\[
                               \|c\|_1\ge K/64.            \tag{Ret.P9}
\]
This counts total Fourier mass, not occupied sites.

To verify (Ret.P3), write \(A^m=\alpha_mI+b_mB\), \(B=A-3I\),
so \(B^2=8I\) and \(\alpha_m^2-8b_m^2=1\). Then
\[
 A^K-I=
 \begin{pmatrix}
  16N^2&8\alpha_mN\\4\alpha_mN&16N^2
 \end{pmatrix}.
\]
Every grid point in (Ret.P3) is fixed by \((A^{\mathsf T})^K\).
The grid has the determinant cardinality and is therefore the whole group.
This also proves \(A^K-I\in2NM_2(\mathbb Z)\).

\hypertarget{the-local-weyl-area-and-the-parity-sign}{%
\subsection{The local Weyl area and the parity sign}\label{the-local-weyl-area-and-the-parity-sign}}

Take integer \(p\) and \(c=\mathcal Dp\).
Put \(z_j=A^jp_j\) for \(0\le j<K\), and \(z_K=A^Kp_0\).
The time-ordered frequencies are \(A^jc_j=z_j-z_{j+1}\).
Their sum is \((I-A^K)p_0\in2N\mathbb Z^2\), so their final
balanced Weyl operator is the identity.
The ordered symplectic area telescopes to
\[
 \sum_{i<j}\omega(A^ic_i,A^jc_j)
   =-\omega(p_0,A^Kp_0)+\sum_j\omega(p_j,Ap_{j+1}).
\]
The first term is divisible by \(2N\). The normalized trace of the
entire word is therefore exactly
\[
 \exp\left(\frac{\pi i}{N}\mathcal S(p)\right),\qquad
              \mathcal S(p)=\sum_j\omega(p_j,Ap_{j+1}).     \tag{Ret.P10}
\]
Since \(\mathcal S(p)=-\sum_j\omega(p_j,c_j)\), (Ret.P8) gives
\[
 |\mathcal S(p)|\le2\|c\|_1^2,\qquad
 \left|e^{\pi i\mathcal S(p)/N}-1\right|
                         \le\frac{2\pi}{N}\|c\|_1^2.       \tag{Ret.P11}
\]
This bound does not grow with the largest frequency \(A^jc_j\):
a complete arithmetic return cancels that growth.

The convention matters. For \(p_0=(1,0)\), \(p_1=(0,1)\), and all
other \(p_j=0\), the local area is 3.
Using the modular inverse of two without \((-1)^{v_1v_2}\)
would give the negative of the small phase in (Ret.P10).
The checker includes odd local areas at every tested period.

\hypertarget{from-the-physical-kick-to-scalar-fourier-coefficients}{%
\subsection{From the physical kick to scalar Fourier coefficients}\label{from-the-physical-kick-to-scalar-fourier-coefficients}}

Let \(k_v(\epsilon)\) be the Fourier coefficients of \(e^{i\epsilon f}\).
They are not the coefficients of \(e^{i\epsilon H_N}\):
Fourier modes within a physical kick need not commute.
For an ordered word of \(\ell\) Weyl frequencies, (Ret.P1) gives
\[
 |\hbox{Weyl phase}-1|
              \le\frac{\pi}{N}\sum_{i<j}\|v_i\|_1\|v_j\|_1.
\]
Summing its exponential series against the absolute Fourier
coefficients proves
\[
 \|e^{i\epsilon H_N}-\operatorname{Op}_N(e^{i\epsilon f})\|
       \le\frac{\pi\delta^2F_1^2}{2N}e^{\delta F_0}.        \tag{Ret.P12}
\]
Both factors have norm at most \(e^{\delta F_0}\).
Replacing all \(K\) interaction-picture kicks therefore costs at most
\(\pi K\delta^2F_1^2e^u/(2N)\) in operator norm.
Commuting the original Fourier modes without this estimate
would not be justified.

The normalized trace of the resulting Fourier product selects
\(\mathcal C_K\). Finite-group Fourier orthogonality in (Ret.P2)
selects exactly \(\mathcal D(\mathbb Z^{2K})\):
the dual kernel equations are \(q_j=A^{\mathsf T}q_{j-1}\) cyclically.
For an integer preimage, (Ret.P10) is the only remaining difference
from the classical coefficient, whose phase is one.

The absolute Fourier mass is at most \(e^u\).
Expanding into Fourier atoms also gives the moment bound
\[
 \sum_{(c_j)}\prod_j|k_{c_j}(\epsilon)|\,\|c\|_1^2
       \le[K\delta F_2+(K\delta F_1)^2]e^u.                 \tag{Ret.P13}
\]
The atom norm bounds the aggregate norm; the right side is the second
moment of the positive exponential majorant. No derivative of an
absolute Fourier coefficient is taken.
Equations (Ret.P11)--(Ret.P13) give the \(N^{-1}\) part of (Ret.P4).

For the other 31 classes, each Taylor atom has norm at most \(M\).
By (Ret.P9), each such term has at least \(L=\lceil K/(64M)\rceil\)
atoms. The absolute mass at degree \(\ell\) is at most
\((K\delta F_0)^\ell/\ell!\), giving \(E_L(u)\).
All classes and Taylor degrees have been included.
For the odd correction, a noncentral Weyl word contributes at most
\(1/(N-1)\) to the difference of normalized traces, and a central word
has zero correction. The atom sum for the actual physical kick is \(e^u\).
This proves Theorem 1.

\hypertarget{smooth-truncation-and-the-time-range}{%
\subsection{Smooth truncation and the time range}\label{smooth-truncation-and-the-time-range}}

Truncate a real smooth \(f\) to \(\|v\|_1\le M\).
The norms \(F_0,F_1,F_2\) stay bounded, and for every \(p\) the Fourier
tail is at most \(C_{f,p}M^{-p}\).
Unitary telescoping and the Lipschitz estimate for real exponentials give
\[
 \left|\frac{\operatorname{Tr}U_{t/\sqrt K,f}^K}{N}
       -\frac{\operatorname{Tr}U_{t/\sqrt K,f_M}^K}{N}\right|
 \le C_{f,p}T\sqrt K\,M^{-p},
\]
\[
 |Z_{K,f}(t/\sqrt K)-Z_{K,f_M}(t/\sqrt K)|
                         \le C_{f,p}T\sqrt K\,M^{-p}.       \tag{Ret.P14}
\]
For the first inequality use Duhamel's formula for self-adjoint
exponentials and the operator norm bound by the Fourier norm.

Choose \(M=\lfloor K^{1/4}\rfloor\).
Then \(u=O_{f,T}(\sqrt K)\) and \(L\ge cK^{3/4}\).
For \(L>u\) the exponential Markov bound gives
\[
 E_L(u)\le(eu/L)^L\le e^{-c_{f,T}K^{3/4}\log K}
\]
for large \(K\). The other term in (Ret.P4) is
\(N^{-1}e^{O(\sqrt K)}O(K)\).
Choose \(p>4q+2\) in (Ret.P14) to prove (Ret.P5).
The odd correction is exponentially small as well.

For a fixed Fourier polynomial and a fixed
\(\delta F_0<1/(128eM)\), instead \(eu/L\le1/2\).
Thus \(E_L(u)\le2^{-L}\) and the other error in (Ret.P4)
also decays exponentially in \(K\). The zero symbol is trivial.
This proves the fixed-small-kick assertion.

For a classical coboundary \(f=g-g\circ A^{\mathsf T}\), (Ret.P2)
telescopes at every periodic point: \(Z_{K,f}=1\).
Theorem 2 correctly retains a coherent first return for these symbols.
No positive-variance claim is made for arbitrary phase-space symbols.

A classical limit theorem is a separate input to one-return asymptotics;
it is not used in the proof of (Ret.P5). Section 7 supplies the missing passage
by a direct coupling to Haar measure and a verified Haar CLT.
The following unsuccessful source leads are retained to identify what
is not being used.
The flow theorem of Cantrell and Sharp,
\href{https://arxiv.org/pdf/1805.05692v1.pdf}{arXiv:1805.05692v1},
was inspected, but its Theorem 1.1 concerns flows and is not used
as a theorem about these discrete cat maps.
Its reference to Coelho--Parry, Section 6, is a source lead,
not a checked proof input. The
\href{https://link.springer.com/article/10.1007/BF02937307}{publisher record}
confirms that the paper discusses closed orbital measures, but the full
section was not accessible in this run and is not cited as verified.

\hypertarget{verification-and-the-remaining-obstruction}{%
\subsection{Verification and the remaining obstruction}\label{verification-and-the-remaining-obstruction}}

\href{https://github.com/ChurilovMaksim/BGS/blob/main/computations/verification/periodic_point_trace_comparison.py}{periodic\_point\_trace\_comparison.py}
checks the determinant, all 32 fractional classes, 960 seeded
rational-preimage configurations, and 1,206 exact integer-word phases
at \(K=6,10,14,18,22,34\).
It includes 603 odd local areas, testing the sign in (Ret.P1).
The physical matrix diagnostic at \(N=35\) uses position, momentum,
and mixed Fourier modes. Double-precision Egorov residuals are
below \(5\cdot10^{-15}\).
It compares the physical trace with all 39,200 classical periodic points.
Differences at \(\epsilon=.08,.04,.02,.01\) are approximately
\(1.82\cdot10^{-3},2.38\cdot10^{-4},3.02\cdot10^{-5},3.79\cdot10^{-6}\).
These are finite diagnostics, not interval bounds or replacements
for the analytic proof.

\begin{verbatim}
OPENBLAS_NUM_THREADS=1 python computations/verification/periodic_point_trace_comparison.py \
  --output build/periodic-point-trace-comparison-precheck.json
\end{verbatim}

The cancellation in (Ret.P10) uses one complete arithmetic period.
Repeating the word introduces interactions between replica assignments.
Two natural extensions have distinct, explicitly demonstrable failures.

First replace the classical period in (Ret.P2) by \(rK\) while keeping the
quantum dimension \(N=b_{K/2}\). For \(f=2\cos(2\pi z_1)\) the exact
quadratic coefficients are
\[
 [\epsilon^2]\frac{\operatorname{Tr}U_{\epsilon,f}^{rK}}N
        =-Kr^2,\qquad
 [\epsilon^2]Z_{rK,f}(\epsilon)=-Kr.                         \tag{Ret.P15}
\]
Here \(Z_{rK,f}\) averages over the actual classical fixed points of
\((A^{\mathsf T})^{rK}\), not the smaller grid \(\mathcal P_K\).
For the quantum coefficient each orbit frequency occurs \(r\) times;
opposite frequencies give \(r^2\) matching pairs.
Within a single period the orbit vectors are distinct up to sign.
Indeed \(0<b_d<N\) for \(0<d<m\), and
\(A^m=\alpha_mI\pmod N\) reduces \(m<d<2m\) to that range.
At \(d=m\) the scalar \(\alpha_m\) is neither \(1\) nor \(-1\) modulo \(N\):
\(2N<\alpha_m<3N\) and the Pell identity excludes
\(\alpha_m=2N+1\) and \(\alpha_m=3N-1\) for the present \(N\ge35\).
The classical period-\(rK\) grid has no such earlier repeated frequency.
Writing \(N'=b_{rK/2}\), a frequency difference has second coordinate
\(2b_d\) modulo \(4N'\); it cannot vanish for \(0<d<rK\).
The only possible multiple of \(N'\) occurs at \(d=rK/2\), where
\(2b_d=2N'\) is still nonzero modulo \(4N'\).
This proves (Ret.P15) by character orthogonality.
The exact checker confirms these two counts at \(K=10\), \(1\le r\le8\).

Second keep the classical period \(K\) and replace its kick by \(r\epsilon\),
which repairs the quadratic discrepancy. The quartic calculation in
\href{https://github.com/ChurilovMaksim/BGS/blob/main/mathematics/fixed_harmonic_return_gas.md}{the fixed-harmonic note} gives, for every
fixed \(r\ge2\) and real \(t\),
\[
 K\left[
 \frac{\operatorname{Tr}U_{t/(r\sqrt K),f}^{rK}}N
                  -Z_{K,f}(t/\sqrt K)\right]
 \longrightarrow
 -\frac{1-r^{-2}}3 S_1t^4e^{-t^2},\qquad
 S_1=\sum_{\ell\ge1}(1-\cos(4\pi\lambda^{-\ell}))>0.         \tag{Ret.P16}
\]
To justify this consequence, use that note's uniform sixth-order
remainder at \(x=r\epsilon=t/\sqrt K\): it is \(O(Kx^6)=O(K^{-2})\).
Its nonzero-aggregate error is \(O(Kx^8)\) plus exponential terms.
Apply (Ret.P4) to the single-period classical comparison; its error at this
scale is exponentially small. The two traces have the same Gaussian
factor, and the displayed quartic difference gives (Ret.P16).
Thus even this repaired extension of the superpolynomial estimate
(Ret.P5) is false: its error generally starts at order \(K^{-1}\).
This is a failure of the proposed comparison, not a refutation of BGS.

Neither this comparison nor the all-replica position-symbol bounds
control the joint traces at \(n\asymp N\) needed for microscopic BGS.
A fixed nonzero classical kick instead requires quantum amplitude
of order \(N\), outside the useful range of (Ret.P4).

\hypertarget{a-centered-time-coupling-closes-the-classical-limit}{%
\subsection{A centered-time coupling closes the classical limit}\label{a-centered-time-coupling-closes-the-classical-limit}}

Write \(T=A^{\mathsf T}\) and use the torus distance induced by the maximum
norm. For a Lipschitz real function let \(\operatorname{Lip}(f)\) refer to
this distance. Haar measure is normalized to mass one.

\textbf{Lemma 3 (coupling).} There is a coupling of a uniform \(q\in\mathcal P_K\)
and a Haar point \(z\in\mathbb T^2\) for which
\[
 \left|\sum_{j=-m}^{m-1} f(T^jz)-\sum_{j=-m}^{m-1}f(T^jq)\right|
                    \le3\operatorname{Lip}(f).             \tag{Ret.P17}
\]
Consequently the Wasserstein distance of order one between the normalized
Birkhoff sums for these two measures is at most
\(3\operatorname{Lip}(f)/\sqrt K\). In particular,
\[
 \left|Z_{K,f}(t/\sqrt K)-
   \int_{\mathbb T^2}e^{it\sum_{j=0}^{K-1}f(T^jz)/\sqrt K}\,dz\right|
       \le \frac{3|t|\operatorname{Lip}(f)}{\sqrt K}.
\]

Partition the torus into centered rectangular cells of side lengths
\(1/(8N)\) and \(1/(4N)\). Couple each point to its cell center. The cells
have equal Haar mass; ties have measure zero. A lift of \(z-q=\eta\)
satisfies \(|\eta_1|\le1/(16N)\) and \(|\eta_2|\le1/(8N)\).
The two spectral projectors of \(T\) have maximum-norm operator norm
\((1+\sqrt2)/2\). Hence
\[
 \|T^j\|_\infty\le(1+\sqrt2)\lambda^{|j|},\qquad
 \sum_{j=-m}^{m-1}\lambda^{|j|}
             =\frac{(\lambda+1)(\lambda^m-1)}{\lambda-1}.
\]
Using \(N=(\lambda^m-\lambda^{-m})/(4\sqrt2)\), the sum of the
distances is at most
\[
 \frac{1+\sqrt2}{8N}\,
       \frac{(\lambda+1)(\lambda^m-1)}{\lambda-1}
 \le \frac{1+\sqrt2}{1-\lambda^{-6}}<3.
\]
This proves (Ret.P17). The centered sum at a periodic point equals its forward
period sum. The centered Haar sum has the forward sum's distribution by
invariance. These two observations are necessary: the corresponding
pointwise estimate for two forward sums need not be bounded uniformly.

For a real smooth mean-zero symbol define
\[
 \sigma_A^2(f)=\sum_{j\in\mathbb Z}\int f(z)f(T^jz)\,dz
       =\sum_{\mathcal O}
             \left|\sum_{v\in\mathcal O}\widehat f_v\right|^2. \tag{Ret.P18}
\]
Here \(\mathcal O\) runs over oriented nonzero \(A\)-orbits, so an orbit and
its negative are counted separately. Absolute convergence follows because
for each pair of nonzero integer frequencies at most one integer iterate
identifies them; the total absolute correlation sum is at most \(F_0^2\).
This is the existing classical orbit-variance formula, not a new claim
of priority; see \href{https://github.com/ChurilovMaksim/BGS/blob/main/mathematics/resonance-cohomology/weighted_wiener_resonance.md}{the earlier derivation}.

\textbf{Theorem 4 (general one-return law).} For every fixed real mean-zero
\(f\) with \(F_4<\infty\), in particular every smooth \(f\),
\[
 \frac{\operatorname{Tr}U_{t/\sqrt K,f}^{K}}N
        \longrightarrow e^{-\sigma_A^2(f)t^2/2},            \tag{Ret.P19}
\]
uniformly on compact real \(t\)-intervals. The odd trace obeys the same law
when \(f\) is even, and \(\sigma_A^2(f)=0\) is allowed.

The external input is Cohen--Conze, arXiv:1304.4556v2, Theorem 2.18:
the Haar CLT for a totally ergodic toral action and a real mean-zero
observable with absolutely convergent Fourier series. Its assumptions
hold here: the eigenvalues \(\lambda^{\pm1}\) imply total ergodicity,
\(F_4<\infty\) gives the required Fourier convergence, and intervals
\(\{0,\ldots,K-1\}\) are a Følner sequence. The exact source and scope are
recorded in \href{https://github.com/ChurilovMaksim/BGS/blob/main/literature/periodic_return_clt_sources.md}{the source check}.
For zero variance the limit also follows directly from the variance
formula and convergence to zero in \(L^2\).
For a smooth symbol combine this input with (Ret.P17) and (Ret.P5).
For finite \(F_4\), truncate at \(M=\lfloor K^{1/4}\rfloor\) instead:
the two tail errors in (Ret.P14) are at most \(T\sqrt K F_4M^{-4}=O(K^{-1/2})\),
while (Ret.P4) still has a superpolynomially small right side, since
\(F_0,F_1,F_2\) are uniformly bounded. This suffices for the same limit.
For compact uniformity, the Haar characteristic functions have Lipschitz
constant at most \(F_0\), by their second-moment bound. A finite mesh,
pointwise convergence, and the two uniform comparisons finish the proof.
No quantitative rate for the external Gaussian approximation is assumed.
For a nonzero mean \(\mu\), apply the statement to \(f-\mu\) and multiply
the quantum trace by \(e^{-i\mu t\sqrt K}\).

\hypertarget{pair-multiplicity-and-a-finite-dimensional-unitary-normal-form}{%
\subsection{Pair multiplicity and a finite-dimensional unitary normal form}\label{pair-multiplicity-and-a-finite-dimensional-unitary-normal-form}}

The next step adapts the averaging mechanism of
\href{https://github.com/ChurilovMaksim/BGS/blob/main/mathematics/pell_smooth_resonance.md}{the earlier odd-prime-period note}, Section 4,
to the present even arithmetic periods. The 32-class lattice makes it
possible to prove the required pair bound directly, without assuming
that different Fourier modes occupy different classical orbits.

Let \(\mathcal S\subset\mathbb Z^2\setminus\{0\}\) be a finite frequency set
with \(\|v\|_1\le M\). Assume \(K>256M\) and put \(L=8M+1\).
Every nonzero key in \((\mathbb Z/N\mathbb Z)^2\) has at most
\[
 B_M=\left[2(2L+1)|\mathcal S|\right]^2                     \tag{Ret.P20}
\]
ordered representations \(A^iv+A^jw\) with \(0\le i,j<K\) and
\(v,w\in\mathcal S\). Equal times are included.

To prove this, fix one representation and compare another. Their difference
is a labeled four-atom relation with total norm at most \(4M\). Aggregate
the atoms at each time into \(c\). Equation (Ret.P9) forces its preimage to be
integer, and (Ret.P8) gives \(\|p\|_1\le8M\). There is a zero coordinate of \(p\).
Split its nonzero coordinates into maximal contiguous runs. Each run
has length at most \(8M\), and its forcing is supported on that run and
the one preceding site. Forcing intervals from different runs are disjoint.
After cutting the time circle at a zero preimage coordinate, each run
gives an exact integer-frequency zero relation by telescoping.
For a run crossing the original time origin use the lifted time indices;
these agree with the original indices modulo \(N\) because \(A^K=I\bmod N\).

Assign all atoms at a nonzero forcing site to its run. At any remaining
occupied site the aggregate is zero; those atoms form a separate zero
group. Each group contains at least two atoms. A nonempty run has
nonzero forcing at both endpoints; a single nonzero integer frequency
cannot form a zero relation.
There are thus either one group of four atoms or two groups of two.
In the latter case each group must match an atom from each representation:
an internal match would make the fixed output key zero.
It follows that both new times are in the union of the cyclic
radius-\(L\) neighborhoods of the two fixed times. This union has at most
\(2(2L+1)\) sites. Counting the two new frequencies proves (Ret.P20).
The argument includes coincident atoms, equal times, canceled aggregates,
and a vanishing entire forcing vector. Equal times in the first
representation do not permit an internal zero group for a nonzero key.

For \(H=\operatorname{Op}_N(f)\) with this Fourier support, define
\[
 H_j=U_0^jHU_0^{-j},\quad
 \bar H=\frac1K\sum_{j=0}^{K-1}H_j,\quad
 c_j=\frac{K-1-2j}{2K},\quad G=\sum_{j=0}^{K-1}c_jH_j .
\]
Both \(G,\bar H\) are self-adjoint; \(\bar H\) commutes with \(U_0\), and
\(G-U_0GU_0^{-1}=H-\bar H\).
Let
\[
 a_d=\max_{v,w\in\mathcal S}
       |e^{2\pi i\omega(v,A^dw)/N}-1|.
\]
Taking the centered representative \(d\) modulo \(K\), using
\(A^K=I\bmod N\), and summing geometric tails from \(|d|=m\) gives
\[
 \sum_{d=0}^{K-1}a_d\le C_A M^2,\qquad
 \sum_{d=0}^{K-1}|c_d|a_d\le C_A M^2/K.                    \tag{Ret.P21}
\]
More explicitly, if \(\ell=m-\min(d,K-d)\), then
\(a_d\le C_A M^2\lambda^{-\ell}\) and
\(|c_d|\le(\ell+1)/K\). The same estimates include \(d=0\).

Use the normalized Hilbert--Schmidt norm
\(\|X\|_{2,N}^2=N^{-1}\operatorname{Tr}(X^*X)\).
The triangle inequality and (Ret.P21) give
\[
 \|[G,H]\|_{\rm op}+\|[H,\bar H]\|_{\rm op}
                       \le C_A F_0^2 M^2/K .
\]
For the remaining commutator,
\[
 [G,\bar H]=\frac1{K^2}\sum_{i<j}(j-i)[H_i,H_j].
\]
Balanced Weyl matrices in distinct residue classes are orthogonal.
The zero output key has zero commutator coefficient.
Cauchy--Schwarz on each nonzero key, followed by (Ret.P20), therefore yields
\[
 \begin{split}
 \|[G,\bar H]\|_{2,N}^2
 &\le\frac{B_MF_0^4}{K^4}
          \sum_{d=1}^{K-1}(K-d)d^2a_d^2\\
 &\le\frac{C_A B_MF_0^4 M^2}{K}.
 \end{split}                                               \tag{Ret.P22}
\]
The last inequality uses \(d(K-d)\le K^2/4\) and \(a_d^2\le2a_d\).
Since \(|\mathcal S|\le(2M+1)^2\), it follows that
\(\|[G,\bar H]\|_{2,N}\le C_A F_0^2M^4/\sqrt K\).
This last bound is a Hilbert--Schmidt bound, not an operator-norm bound.

For self-adjoint \(X,Y\), integration over a time triangle proves
\[
 \|e^{i\epsilon X}e^{i\epsilon Y}-e^{i\epsilon(X+Y)}\|
                      \le\frac{\epsilon^2}{2}\|[X,Y]\|
\]
in either of these unitarily invariant norms, for every real \(\epsilon\).
Apply it first to \(-G,H\) and then to \(H-G,U_0GU_0^{-1}\).
The second commutator is \([H-G,\bar H]\). Hence
\[
 \left\|e^{-i\epsilon G}U_{\epsilon,f}e^{i\epsilon G}
                          -e^{i\epsilon\bar H}U_0\right\|_{2,N}
 \le C_A F_0^2\epsilon^2
           \left(\frac{M^2}{K}+\frac{M^4}{\sqrt K}\right).    \tag{Ret.P23}
\]
No smallness assumption on \(\epsilon\|G\|\) has been used.
Constants in \(f\) can be removed exactly; they do not affect commutators.

\hypertarget{every-replica-at-the-critical-scale}{%
\subsection{Every replica at the critical scale}\label{every-replica-at-the-critical-scale}}

\textbf{Theorem 5 (uniform replica comparison).} For every fixed real symbol
\(f\) with \(F_4<\infty\) and every \(T<\infty\),
\[
 \sup_{\substack{r\ge1\\|t|\le T}}
 \left|\frac{\operatorname{Tr}U_{t/(r\sqrt K),f}^{rK}}N
          -\frac{\operatorname{Tr}U_{t/\sqrt K,f}^{K}}N\right|
                  =O_{f,T}(K^{-1/2}).                      \tag{Ret.P24}
\]
The result also holds for the normalized odd traces of even symbols.
It concerns integer \(r\), including \(r\) depending on \(K\).

We first remove the bandwidth loss in (Ret.P23). Subtract the mean throughout
this argument; restoring it multiplies both traces by the same phase.
Split the nonzero Fourier frequencies into the shells
\(0<\|v\|_1\le1\) and \(2^{a-1}<\|v\|_1\le2^a\) for \(a\ge1\).
Let \(f_a\) be the corresponding symbol and
\(\mu_a=\sum|\widehat{f_a}|\). The zero mode is excluded.
For two shells of maximal radius \(M\), with \(K>256M\), (Ret.P20) and
Weyl orthogonality applied to
\([G_a,\bar H_b]=K^{-1}\sum_{i,j}c_i[H_{a,i},H_{b,j}]\) give
\[
 \|[G_a,\bar H_b]\|_{2,N}^2
 \le\frac{B_M\mu_a^2\mu_b^2}{K^2}
              \sum_{i,j}c_i^2a_{j-i}^2
 \le \frac{C_A B_M\mu_a^2\mu_b^2M^2}{K}.
\]
Here \(\sum_i c_i^2\le K/4\) and \(\sum_da_d^2\le2\sum_da_d\le C_AM^2\).
The same bound includes equal-time pairs. Thus the bilinear norm is at
most \(C_A\mu_a\mu_b\max(2^a,2^b)^4/\sqrt K\).

Take the largest dyadic \(R<K/512\) and sum over shells of radius at most
\(R\). There is no logarithmic loss, since
\[
 \sum_{a,b}\mu_a\mu_b\max(2^a,2^b)^4
 \le2\left(\sum_a\mu_a\right)
       \left(\sum_b2^{4b}\mu_b\right)\le C F_0F_4.
\]
The omitted commutators are
\([G_{>R},\bar H]+[G_{\le R},\bar H_{>R}]\).
Their operator norm is at most \(CKF_0\sum_{\|v\|_1>R}|\widehat f_v| \le CK^{-3}F_0F_4\). For the finitely many smaller \(K\), enlarge the constant.
The full operator-norm bounds also follow by summing (Ret.P21) first for
each frequency pair and then against its Fourier coefficients. Therefore
\[
 \begin{split}
 \|[G,\bar H]\|_{2,N}&\le C_A F_0F_4/\sqrt K,\\
 \|[G,H]\|_{\rm op}+\|[H,\bar H]\|_{\rm op}
                         &\le C_A F_1^2/K,\\
 \left\|e^{-i\epsilon G}U_{\epsilon,f}e^{i\epsilon G}
                         -e^{i\epsilon\bar H}U_0\right\|_{2,N}
 &\le C_A\epsilon^2
                   (F_1^2/K+F_0F_4/\sqrt K).
 \end{split}                                               \tag{Ret.P24a}
\]
All Fourier passages are norm-convergent. No truncation of the actual
quantum trace is required here.

Telescope \(rK\) unitary factors in (Ret.P24a), with
\(\epsilon=t/(r\sqrt K)\). The reference product is exactly
\(e^{it\sqrt K\bar H}\); its normalized trace differs from the physical
trace by at most
\(C_AT^2r^{-1}(F_1^2/K+F_0F_4/\sqrt K)\).
Comparing with \(r=1\) proves (Ret.P24). The gauge depends on \(\epsilon\)
but disappears from each trace.
For even symbols all gauges preserve parity; the restriction of a
normalized full Hilbert--Schmidt bound to the odd space costs at most
\(\sqrt{2N/(N-1)}\).

Combining Theorems 4 and 5 gives, for every fixed real mean-zero \(f\)
with \(F_4<\infty\),
\[
 \sup_{\substack{r\ge1\\|t|\le T}}
 \left|\frac{\operatorname{Tr}U_{t/(r\sqrt K),f}^{rK}}N
                    -e^{-\sigma_A^2(f)t^2/2}\right|\to0.    \tag{Ret.P25}
\]
In particular, if \(\epsilon_K\sqrt K\to s\in\mathbb R\), the global
empirical eigenphase distribution of \(U_{\epsilon_K,f}^K\) converges to
the wrapped Gaussian of variance \(s^2\sigma_A^2(f)\). Indeed each fixed
integer moment tends to \(e^{-r^2s^2\sigma_A^2(f)/2}\); negative moments
are conjugates, and trigonometric polynomials are dense on the circle.
For a nonzero mean \(\mu\), center the return operator by
\(e^{-iK\epsilon_K\mu}\). The odd-block conclusion requires even \(f\).

The \(K^{-1/2}\) estimate is for the comparison of quantum returns;
the Gaussian limit itself is qualitative here.
The nonzero \(K^{-1}\) correction in (Ret.P16) is consistent with (Ret.P24).
This does not extend the large-variance decay estimate (Ret.A1) to all
two-dimensional smooth symbols. The separate
\href{https://github.com/ChurilovMaksim/BGS/blob/main/mathematics/phase_space_return_gas.md}{analytic phase-space argument} covers
general analytic symbols in that regime using a vector-lattice gas
estimate. Neither argument determines local level
spacings, or traces at the Heisenberg scale for one fixed perturbed
classical system.

\hypertarget{the-zero-variance-case-and-its-next-scale}{%
\subsection{The zero-variance case and its next scale}\label{the-zero-variance-case-and-its-next-scale}}

\textbf{Theorem 6 (smooth degeneracy).} For real mean-zero smooth \(f\),
\(\sigma_A^2(f)=0\) if and only if \(f=g-g\circ T\) for a real smooth \(g\).
If \(f\) is even, \(g\) can be chosen even.
This is a smooth Fourier-coboundary criterion, a classical mechanism
already used in \href{https://github.com/ChurilovMaksim/BGS/blob/main/mathematics/pell_smooth_resonance.md}{the earlier smooth note}.
Here is the regularity argument for the current covector matrix \(A\).

Choose real eigen-coordinates \(\xi,\zeta\) with
\(A:(\xi,\zeta)\mapsto(\lambda\xi,\lambda^{-1}\zeta)\) and let
\(|v|_*=\max(|\xi(v)|,|\zeta(v)|)\). No nonzero integer vector lies on
either eigenline. On each oriented orbit select a representative \(v_0\)
minimizing this norm, paired with \(-v_0\) on the negative orbit.
Minimality implies
\(\lambda^{-1}\le|\xi(v_0)/\zeta(v_0)|\le\lambda\).
Writing \(v_j=A^jv_0\), one obtains
\(\lambda^{|j|-1}|v_0|_*\le|v_j|_*\le\lambda^{|j|}|v_0|_*\).

If the orbit sum of \(f\) vanishes, define its transfer coefficients by
\[
 \widehat g_{v_j}=
 \begin{cases}
   \sum_{\ell\le j}\widehat f_{v_\ell},&j<0,\\
   -\sum_{\ell>j}\widehat f_{v_\ell},&j\ge0 .
 \end{cases}
\]
For every \(s>0\), summing the geometric weights on each half-orbit gives
\(\sum_j|v_j|_*^s|\widehat g_{v_j}|  \le C_s\sum_j|v_j|_*^s|\widehat f_{v_j}|\).
For example, on the positive half use
\(|v_j|_*^s\le\lambda^s\lambda^{-s(\ell-j)}|v_\ell|_*^s\)
for \(\ell>j\ge0\), and interchange absolutely convergent sums.
The negative half is identical with time reversed.
Summing over all orbits proves smoothness of \(g\) by equivalence of norms.
Its coefficient difference is exactly \(f\), including at \(j=0\) because
the orbit sum vanishes. Reality and parity are preserved by the construction.
Conversely, the Fourier differences of a smooth coboundary telescope
on every orbit. Formula (Ret.P18) proves the equivalence.

Let \(X=\operatorname{Op}_N(g)\), \(Y=\operatorname{Op}_N(g\circ T)\).
The elementary two-factor Duhamel estimate gives
\[
 \left\|e^{-i\epsilon X}U_{\epsilon,f}e^{i\epsilon X}-U_0\right\|_{\rm op}
 \le\frac{\epsilon^2}{2}\|[X,Y]\|_{\rm op}
 \le\frac{\pi\epsilon^2}{N}F_1(g)F_1(g\circ T).              \tag{Ret.P26}
\]
Consequently, at \(\epsilon=t/(r\sqrt K)\) the normalized \(rK\)-return
trace differs from one by at most \(C_g t^2/(rN)\). For \(r=1\) even the
operator \(U_{\epsilon,f}^K\) is within \(C_gt^2/N\) of the identity.
This is stronger than the degenerate case of the qualitative limit.

The next effect is described by the fixed smooth symbol
\[
 h_g=\frac{\{g,g\circ T\}}{4\pi},\qquad
 \Gamma_A(g)=\sigma_A^2(h_g),\qquad
 \{a,b\}=\partial_1a\,\partial_2b-\partial_2a\,\partial_1b .
\]
It is real and mean zero; it is even if \(g\) is even.
For self-adjoint \(X,Y\), put \(Q=[X,Y]\). The uniform unitary loop bound is
\[
 \left\|e^{-i\epsilon X}e^{i\epsilon(X-Y)}e^{i\epsilon Y}
                             -e^{-\epsilon^2Q/2}\right\|_{\rm op}
 \le|\epsilon|^3\left(\tfrac12\|[X,Q]\|
                              +\tfrac16\|[Y,Q]\|\right).   \tag{Ret.P27}
\]
Indeed its logarithmic derivative is
\(-\int_0^u\operatorname{Ad}_{e^{-iuX}e^{iv(X-Y)}}Q\,dv\).
Compare it with \(-uQ\), the logarithmic derivative of
\(e^{-u^2Q/2}\), using unitary conjugation, and integrate over \(u\).
The bound is valid for every real \(\epsilon\); it does not require
\(\epsilon\|X\|\) to be small.

The Weyl formula and \(|\sin u-u|\le|u|^3/6\) imply
\([X,Y]=-2i\operatorname{Op}_N(h_g)/N+O_g(N^{-3})\).
Both double commutators in (Ret.P27) are \(O_g(N^{-2})\) by two applications
of \(|2\sin(\pi\omega/N)|\le2\pi|\omega|/N\) to the Fourier series.
All weighted sums converge for smooth \(g\), with no frequency exclusion.
Thus
\[
 \left\|e^{-i\epsilon X}U_{\epsilon,g-g\circ T}e^{i\epsilon X}
                             -U_{\epsilon^2/N,h_g}\right\|_{\rm op}
       \le C_g\left(\frac{|\epsilon|^3}{N^2}
                                +\frac{\epsilon^2}{N^3}\right). \tag{Ret.P28}
\]
This is the same local loop mechanism established in the earlier note;
the present consequence concerns the even-period sequence and all replicas.

\textbf{Theorem 7 (uniform second-order return law).} For every fixed real
smooth \(g\) and every finite \(T_0\),
\[
 \sup_{\substack{r\ge1,\ 0\le t\le T_0\\
                 \epsilon=\pm\sqrt{tN/(r\sqrt K)}}}
 \left|\frac{\operatorname{Tr}U_{\epsilon,g-g\circ T}^{rK}}N
                            -e^{-\Gamma_A(g)t^2/2}\right|\to0. \tag{Ret.P29}
\]
Here \(r\) ranges over positive integers. The same holds on the odd block
for even \(g\).
After \(rK\) steps the error in (Ret.P28) is at most
\[
 C_{g,T_0}\left(\frac{K^{1/4}}{\sqrt{rN}}
                                  +\frac{\sqrt K}{N^2}\right).
\]
The effective kick is exactly \(t/(r\sqrt K)\), so (Ret.P25) applies to
the fixed symbol \(h_g\), uniformly in \(r\). This proves (Ret.P29).
No positivity of \(\Gamma_A(g)\) for every nonconstant \(g\) is asserted.

For example, take \(g=2\cos(2\pi z_1)\) and \(v=(1,0)\). Then
\[
 h_g=2\pi\bigl(e_{v-Av}+e_{-v+Av}-e_{v+Av}-e_{-v-Av}\bigr),
 \qquad \Gamma_A(g)=16\pi^2.                               \tag{Ret.P30}
\]
The two frequency pairs occupy different signed \(A\)-orbits:
\(\omega(v-Av,A(v-Av))=-8\) and
\(\omega(v+Av,A(v+Av))=16\). These invariants differ and are nonzero.
The limiting second-order trace is therefore \(e^{-8\pi^2t^2}\).
The quantum amplitude in (Ret.P29) can grow, but its classical strength
\(|\epsilon|/N\) still tends to zero uniformly in the indicated range.

\hypertarget{joint-traces-of-physical-return-operators}{%
\subsection{Joint traces of physical return operators}\label{joint-traces-of-physical-return-operators}}

The normal form also controls the return operator itself. For a real
symbol with \(F_4<\infty\), let \(X_K(f)=\sqrt K\,\bar H_f\).
For fixed bounded \(t\) and every positive integer \(r\),
\[
 \left\|U_{t/(r\sqrt K),f}^{rK}-e^{itX_K(f)}\right\|_{2,N}
 \le \frac{C_A t^2}{r}
                   \left(\frac{F_1^2}{K}+\frac{F_0F_4}{\sqrt K}\right).
                                                               \tag{Ret.P31}
\]
To remove the gauge from (Ret.P24a), apply the integrated commutator bound
\[
 \|e^{i\alpha G}e^{i\beta\bar H}e^{-i\alpha G}
                                  -e^{i\beta\bar H}\|_{2,N}
                     \le|\alpha\beta|\|[G,\bar H]\|_{2,N}
\]
with \(\alpha=t/(r\sqrt K)\) and \(\beta=t\sqrt K\).
Together with the telescoped normal form, this proves (Ret.P31).
It does not require the gauge itself to be close to the identity.

For two fixed real symbols \(f,h\), the same ordered-pair argument as in
Section 9 gives
\[
 \|[X_K(f),X_K(h)]\|_{2,N}
 \le\frac{C_A}{\sqrt K}
              \bigl(F_0(f)F_4(h)+F_4(f)F_0(h)\bigr).        \tag{Ret.P32}
\]
For two low-frequency shells the coefficient in
\([\bar H_f,\bar H_h]\) is \(K^{-2}\) at each ordered time pair.
Thus its squared norm is bounded by
\(B_M\mu_a^2\nu_b^2K^{-4}\sum_{i,j}a_{j-i}^2 \le C_AB_M\mu_a^2\nu_b^2M^2/K^3\).
Multiply the norm by \(K\) for the two factors \(X_K\), and sum the
dyadic weights. The high-frequency remainder before multiplication
by \(K\) is \(O(F_0F_{\rm tail})\), hence \(O(K^{-4})\) with the
corresponding weighted norms. Constants can again be adjusted for
bounded \(K\). This proves (Ret.P32), including constants after their removal.

\textbf{Theorem 8 (joint critical return law).} Fix real mean-zero
\(f_1,\ldots,f_J\) with finite fourth Fourier moments and real
\(t_1,\ldots,t_J\). Write
\[
 \Sigma_{ab}=\sum_{\mathcal O}
       \left(\sum_{v\in\mathcal O}\widehat{f_a}_v\right)
       \overline{\left(\sum_{v\in\mathcal O}\widehat{f_b}_v\right)} .
\]
This real symmetric covariance is the polarization of (Ret.P18).
Then, uniformly over positive integers \(r_1,\ldots,r_J\),
\[
 \frac1N\operatorname{Tr}
       \left(\prod_{j=1}^J
               U_{t_j/(r_j\sqrt K),f_j}^{\,r_jK}\right)
 \longrightarrow
       \exp\left(-\tfrac12\sum_{a,b=1}^Jt_at_b\Sigma_{ab}\right).
                                                               \tag{Ret.P33}
\]
The product has the displayed order; every fixed order has the same limit.
For the odd-block statement every \(f_j\) must be even.

Equation (Ret.P31) replaces the physical factors by \(e^{it_jX_K(f_j)}\)
with a total normalized Hilbert--Schmidt error \(O(K^{-1/2})\).
Two-factor Duhamel, iterated over this fixed number of factors, gives
\[
 \left\|\prod_j e^{it_jX_K(f_j)}
          -e^{iX_K(\sum_jt_jf_j)}\right\|_{2,N}
 \le\tfrac12\sum_{a<b}|t_at_b|
                       \|[X_K(f_a),X_K(f_b)]\|_{2,N}.
\]
Apply (Ret.P32), then (Ret.P31) with \(r=1,t=1\) to the fixed symbol
\(\sum_jt_jf_j\), and finally Theorem 4. Its variance is the quadratic
form in (Ret.P33). This proves the joint law without asserting convergence
of unbounded moments from characteristic functions alone.

These are joint traces of a fixed finite family of arithmetic return
operators with the explicitly scaled amplitudes. They are not the
Heisenberg-scale joint traces of a fixed classical perturbation.
For a fixed amplitude \(\epsilon=t/\sqrt K\), the direct normal-form
error after \(rK\) steps grows as \(O_f(rt^2/\sqrt K)\).
Thus its useful fixed-amplitude time range is \(r=o(\sqrt K)\),
whereas Heisenberg times require \(r\) of order \(N/K\).

\hypertarget{verification-scope-of-the-continuation}{%
\subsection{Verification scope of the continuation}\label{verification-scope-of-the-continuation}}

The same verifier now includes seven exact centered-distance bounds,
through \(K=262\), and exhausts 1,122,112 ordered two-atom representations
at \(K=6,10,14,34,262\), including equal-time pairs.
The \(K=262\) case contains 1,098,304 representations and satisfies
the hypothesis \(K>256M\) with \(M=1\). The observed maximum nonzero-key
multiplicity is 4 in every case; the proved bound is 23,104.
The smaller cases are explicitly marked outside that theorem's range.
The discrepancy between this conservative universal bound and these
observations is not used to strengthen the theorem.

At \(N=35\), physical matrices check the exact averaging identity,
five unitary normal-form inequalities, ten replica comparisons both
with and without the gauge, and six joint physical-product comparisons.
For the coboundary example, the generated Fourier coefficients and
the orbit invariants in (Ret.P30) are checked exactly. Nine unitary loops,
including negative and non-small amplitudes, and eight second-order
replica comparisons are tested in double precision against their finite
Duhamel bounds. No interval enclosure or numerical proof of a limit
is claimed.

The independent internal review checked the 32-class and balanced-sign
argument, the ordered-pair decomposition including equal times, the
dyadic \(F_4\) estimate, and the smooth transfer and second-order signs.
The initial P4--P5 snapshot had SHA-256
6324e620fad6a8cb13888b80738de3d0153d6a1b3f3f74bf1897019717247fb0.
The Section 10 audit referenced snapshot
6e79e11455a9b91226eda1b66cd01079ddf9344d1e18a43f01ab6e6e09b9f1f2.
Subsequent amendments make the integer-replica quantifier explicit and
add Section 11; that extension is not represented as independently audited.
The source check documents the external classical CLT separately.
Generated results remain in the reproducible build output. These proofs are
included in the canonical manuscript through a generated TeX chapter; the
analytic limit theorems are not Lean-formalized.
